% ============================================================
% 第11章 生物大分子、合成高分子与有机合成
% ============================================================
\chapter{生物大分子、合成高分子与有机合成}

前两章学习了单个有机分子的结构和官能团反应。本章把这些知识用于两类新对象：一类是糖类、油脂、蛋白质、核酸等与生命有关的物质；另一类是由许多小分子连接而成的高分子材料。最后再反过来思考：给定一种目标物，怎样利用已知反应安排一条合成路线。学习中始终追踪三件事：碳骨架怎样变化，官能团在哪里变化，连接处断开或形成了什么键。

\section{糖类与油脂：从结构认识营养物质}

\subsection{糖类的分类依据是能否水解}
葡萄糖能直接被人体利用，淀粉却要经过消化才能提供葡萄糖。两者的区别不在于“甜不甜”，而在于分子能否被水拆成更小的糖。糖类是多羟基醛、多羟基酮及能水解生成它们的物质；并非所有糖都有甜味，也并非都符合“碳的水合物”的形式。

以葡萄糖为起点。在高中常用的链状结构模型中，葡萄糖写成 \ce{CH2OH(CHOH)4CHO}：一端有醛基，另外五个碳上各有羟基。因此它能与水形成较多氢键、易溶于水；也能被银氨溶液或碱性新制 \ce{Cu(OH)2} 氧化。它已经不能水解成更简单的糖，属于\textbf{单糖}。实际水溶液中葡萄糖主要以环状结构存在，高中讨论上述性质时使用链状模型，不把链状式理解为溶液中的唯一结构。果糖的分子式也是 \ce{C6H12O6}，链状模型为多羟基酮，与葡萄糖互为同分异构体。

两个单糖单元通过连接形成\textbf{二糖}，水解时重新得到两个单糖分子。蔗糖和麦芽糖都为 \ce{C12H22O11}，但连接方式不同：
\[
\ce{C12H22O11 + H2O ->[\text{酸或酶}] C6H12O6 + C6H12O6}
\]
蔗糖的产物是葡萄糖和果糖；麦芽糖的产物只有葡萄糖，1 mol 麦芽糖对应 2 mol 葡萄糖。水解前，麦芽糖有还原性，蔗糖没有，因此不能把“含糖”直接等同于“能发生银镜反应”。果糖虽没有醛基，但在这些碱性检验条件下可发生结构转化而表现还原性；仅凭银镜反应不能唯一鉴定葡萄糖。

许多单糖单元连接形成\textbf{多糖}。淀粉和纤维素都可写作 \ce{(C6H10O5)_{$n$}}，水解的最终产物都是葡萄糖：
\[
\ce{(C6H10O5)_{$n$} + $n$H2O ->[\ce{H+}][\text{△}] $n$C6H12O6}
\]
这里采用教材的重复单元通式，忽略长链端基，水解按该通式计量；后文对有限条直链逐个数连接时会保留端基，二者不能混用。它们的连接结构不同，聚合度 $n$ 也不确定，不能仅凭这一通式认为两者分子式相同、互为同分异构体。它们属于天然高分子，葡萄糖和二糖不属于高分子。

\tu{\begin{tikzpicture}[font=\small,>=Stealth]
\foreach \x[count=\i from 0] in {0,1.8,3.6,5.4}{\node[draw,rounded corners=2pt,minimum width=1.3cm,minimum height=.65cm] (a\i) at (\x,0) {糖单元};}
\draw (a0)--(a1) (a1)--(a2) (a2)--(a3);
\foreach \x in {.9,2.7,4.5}{\draw[red!70!black,dashed] (\x,-.5)--(\x,.5);}
\node at (-1.2,0) {$\cdots$};\node at (6.6,0) {$\cdots$};
\draw[->] (7.2,0)--node[above]{水解}node[below]{加水、断连接}(9.1,0);
\node[align=center] at (10.4,0) {独立的\ce{C6H12O6}\\分子};
\node[align=center,font=\scriptsize] at (2.7,-1.05) {每个方块表示一个葡萄糖单元，连线表示单元之间的连接；\\这不是淀粉或纤维素的详细结构式。};
\end{tikzpicture}}{读多糖水解图：先找重复单元，再看连接处的变化\label{fig:11-sugar}}

\begin{liti}[由水解关系计算]
16.2 g 纯淀粉完全水解，理论上生成多少克葡萄糖？若测得葡萄糖为 14.4 g，能否仅据此断定淀粉只水解了 80\%？
\tcblower
\jie 一个 \ce{C6H10O5} 单元的式量为 162，水解后对应一个式量为 180 的葡萄糖分子。理论质量为 $16.2\times180/162=18.0$ g，增重来自水。$14.4/18.0=80\%$ 是葡萄糖的实际获得量占理论量的比例；若有取样、分离损失或葡萄糖继续反应，不能把它直接当作淀粉的水解率。只有题目说明这些影响可忽略，才可这样计算。
\end{liti}

\subsection{从实验现象判断水解程度}
淀粉遇碘显蓝色，这是检验淀粉的反应，不能检验葡萄糖。检验水解产物则利用其还原性。两个检验回答不同的问题，应该\textbf{另取样品，分别操作}，避免前一种试剂干扰后一种检验。

\tu{\begin{tikzpicture}[font=\small,>=Stealth]
\node[draw,minimum height=.65cm] (s) at (0,0) {淀粉酸性水解液};
\node[draw,align=center,text width=4.2cm] (a) at (5,1) {冷却后加碘水\\看是否仍有淀粉};
\node[draw,align=center,text width=4.2cm] (b) at (5,-1.1) {先用 NaOH 调至碱性\\再加新制\ce{Cu(OH)2}并加热};
\draw[->] (s.east)--(2,0)|-(a.west);\draw[->] (2,0)|-(b.west);
\node[align=left,anchor=west] at (8,1) {蓝色：淀粉仍存在};
\node[align=left,anchor=west] at (8,-1.1) {砖红色沉淀：\\有还原性水解产物};
\end{tikzpicture}}{淀粉水解的分样检验流程\label{fig:11-starchtest}}

为什么第二份必须先加碱？水解用的酸会消耗 \ce{OH-}，还会溶解 \ce{Cu(OH)2}，破坏醛基检验所需的碱性条件。只看“不变蓝”也不够：必须确认原样确含淀粉，且水解产物的检验有效，才能排除样品过稀、试剂失效等原因。

\begin{center}\small
\begin{tabular}{lll}
\toprule
加碘水 & 水解产物的还原性检验 & 在上述实验条件有效时的判断\\
\midrule
变蓝 & 阳性 & 既有淀粉又有水解产物，部分水解\\
不变蓝 & 阳性 & 淀粉已水解完全\\
变蓝 & 阴性 & 未检出水解产物，按高中模型判断未水解\\
\bottomrule
\end{tabular}\end{center}

\subsection{油脂是三酯，不是高分子}
油脂由甘油（丙三醇，\ce{HOCH2CH(OH)CH2OH}）和高级脂肪酸形成。甘油有三个羟基，每个都可与一分子羧酸酯化，所以一个油脂分子有\textbf{三个酯基}。三个脂肪酸烃基可以相同，也可以不同。天然油脂通常是混合物，其分子虽比乙酸乙酯大，却没有高分子那样大量重复的结构单元，不能归为高分子。

\tu{\small
\begin{tabular}{ccc}
\(\begin{array}{l}\ce{CH2-O-CO-R1}\\[-1pt]\vert\\[-1pt]\ce{CH-O-CO-R2}\\[-1pt]\vert\\[-1pt]\ce{CH2-O-CO-R3}\end{array}\)
&\(\xrightarrow[\text{加热}]{\text{足量 NaOH 溶液}}\)&
\(\begin{array}{c}\ce{HOCH2CH(OH)CH2OH}\\+\\\ce{R1COONa + R2COONa + R3COONa}\end{array}\)\\
油脂（每条横向支链有一个酯基）&&甘油与高级脂肪酸钠\\
\end{tabular}}{油脂的结构与皂化产物（R$_1$、R$_2$、R$_3$ 为烃基）\label{fig:11-fat}}

油脂在酸性条件下水解得到甘油和高级脂肪酸，在碱性条件下水解得到甘油和高级脂肪酸盐，后者叫\textbf{皂化反应}。例如硬脂酸甘油酯：
\[
\ce{(C17H35COO)3C3H5 + 3NaOH ->[\text{△}] 3C17H35COONa + C3H5(OH)3}
\]
烃基中含碳碳双键较多的油脂通常熔点较低，常温下呈液态，称油；饱和程度较高的通常熔点较高，常温下呈固态，称脂肪。催化加氢减少双键，可使油硬化；这一步改变的是碳碳双键，不是酯基。判断油脂的反应时必须同时看这两类结构。

\begin{liti}[把两个反应位置分开]
某油脂的每个分子中有 3 个酯基和 3 个碳碳双键。1 mol 该油脂完全皂化消耗多少摩尔 NaOH？完全氢化消耗多少摩尔 \ce{H2}？氢化后还能否皂化？
\tcblower
\jie 皂化在酯基处发生，每个酯基消耗 1 mol NaOH，共 3 mol；氢化在双键处发生，每个双键消耗 1 mol \ce{H2}，共 3 mol。氢化没有破坏酯基，所以产物仍能皂化，且每摩尔仍消耗 3 mol NaOH。
\end{liti}

\section{氨基酸、蛋白质与核酸}

\subsection{氨基酸为什么具有两性}
氨基酸分子中同时含有氨基和羧基。高中常见的 $\alpha$-氨基酸，其两个官能团连在同一个碳上，通式为 \ce{H2N-CH(R)-COOH}。R 为不同侧链：R 为 H 是甘氨酸，为 \ce{CH3} 是丙氨酸。酸碱性不是由名称决定，而是由两个官能团分别与什么反应决定：羧基可失去 \ce{H+}，氨基上的 N 可结合 \ce{H+}，因此能与酸、碱分别反应，是\textbf{两性化合物}。
\[
\ce{H2NCH2COOH + HCl -> [H3NCH2COOH]Cl}
\]
\[
\ce{H2NCH2COOH + NaOH -> H2NCH2COONa + H2O}
\]
在水中还可存在带正、负电荷的两性离子 \ce{H3N+CH2COO-}。高中一般根据题目给定的酸碱条件判断存在形式，不要求用大学酸碱理论计算各形式比例。

\subsection{从两个氨基酸到一条多肽链}
氨基酸连接时，一个分子的羧基脱去 OH，另一分子的氨基脱去 H，生成水，同时在羰基碳与氮之间形成新键。形成的 \ce{-CO-NH-} 叫\textbf{肽键}，它是一种酰胺结构。

\tu{\begin{tikzpicture}[font=\small,>=Stealth]
\node (a) at (0,0) {\ce{H2NCH(R1)CO}};\node (b) at (2.1,0) {OH};\draw (a)--(b);
\node at (2.65,0) {$+$};\node (c) at (3.15,0) {H};\node (d) at (4.6,0) {\ce{NHCH(R2)COOH}};\draw (c)--(d);
\foreach \x in {1.45,3.55}{\draw[red!70!black,dashed] (\x,-.35)--(\x,.35);}
\draw[->] (6.4,0)--(7.4,0);
\node (q1) at (8.3,0) {\ce{H2NCH(R1)}};\node (qc) at (9.7,0) {C};\node (qo) at (9.7,.85) {O};
\node (qn) at (10.7,0) {NH};\node (q2) at (12.3,0) {\ce{CH(R2)COOH}};
\draw (q1)--(qc) (qn)--(q2);\draw[blue!70!black,thick] (qc)--(qn);
\draw ([xshift=-1.5pt]qc.north)--([xshift=-1.5pt]qo.south) ([xshift=1.5pt]qc.north)--([xshift=1.5pt]qo.south);
\node[font=\scriptsize] at (10.2,-.55) {新形成的 C—N 键};\node at (11.1,-1) {$+\ \ce{H2O}$};
\node[font=\scriptsize] at (2.65,-.62) {羧基脱 OH，氨基脱 H};
\end{tikzpicture}}{肽键的形成：先找脱去的原子，再找新连接处（结构变化示意）\label{fig:11-peptide}}

两个氨基酸缩合得到二肽；继续在两端连接可得多肽。若 $m$ 个氨基酸只形成一条\textbf{不成环、不交联的直链}，有 $m-1$ 个肽键，脱去 $m-1$ 个水分子。蛋白质由氨基酸残基连接成长链，并形成特定空间结构；生理功能既与组成和连接次序有关，也与空间结构有关。肽键可在酸、碱或酶作用下水解，充分水解的最终产物是氨基酸（强酸、强碱中实际存在相应的盐）。

\begin{liti}[次序与守恒]
取 1 分子甘氨酸（\ce{C2H5NO2}）和 1 分子丙氨酸（\ce{C3H7NO2}），每次各用一个、只形成直链二肽，可得几种连接次序？二肽分子式是否相同？
\tcblower
\jie 可以是甘氨酸的羧基连接丙氨酸的氨基，也可以反过来，得到
\ce{H2NCH2CONHCH(CH3)COOH} 与 \ce{H2NCH(CH3)CONHCH2COOH}，共 2 种。均脱去一个水，分子式为 \ce{C5H10N2O3}，互为同分异构体。若误把它们写成“甘氨酸二聚体、丙氨酸二聚体”，便违背了每次各用一个的限制。
\end{liti}

\subsection{盐析、变性与检验}
浓的硫酸铵等盐溶液会降低某些蛋白质的溶解度，使它们从溶液中析出；加水稀释后又可溶解，称\textbf{盐析}，常用于分离蛋白质。盐析没有破坏蛋白质原有的结构，通常可逆。

加热、强酸强碱、乙醇、甲醛或某些重金属盐等可改变蛋白质的结构，使其失去原有生理活性，称\textbf{变性}，高中通常按不可逆变化处理。变性并不等于把全部肽键水解成氨基酸：煮熟鸡蛋仍含蛋白质。豆腐凝固涉及盐离子使蛋白质胶体聚集等因素，不宜只凭“加入盐”就把任何凝固一律判为可逆盐析。

含有相应芳香族结构的蛋白质遇浓硝酸可显黄色；灼烧蛋白质有烧焦羽毛的气味，可用于区分蚕丝、羊毛与棉纤维等。酶是生物催化剂，绝大多数酶是蛋白质，具有高效性、专一性，需要适宜温度和酸碱条件；不能说酶全是蛋白质，也不能说升温一定提高酶的催化效果。

\subsection{核酸：骨架的连接与链间的配对}
核酸也是生物大分子，但它的基本单元是\textbf{核苷酸}，不是氨基酸。每个核苷酸由磷酸、戊糖和含氮碱基三部分组成。戊糖种类和部分碱基不同，使核酸分为 DNA 与 RNA。高中重点认识组成、结构层次和氢键的作用，不要求推导核酸合成的详细机理。

\tu{\begin{tikzpicture}[font=\small,>=Stealth]
\foreach \i in {0,1,2}{
\node[draw,circle,inner sep=3pt] (p\i) at (\i*3.1,0) {磷酸};
\node[draw,regular polygon,regular polygon sides=5,inner sep=2pt] (s\i) at (\i*3.1+1.4,0) {戊糖};
\node[draw] (b\i) at (\i*3.1+1.4,1.1) {碱基};
\draw (p\i)--(s\i)--(b\i);}
\draw (s0)--(p1) (s1)--(p2);
\draw[decorate,decoration={brace,mirror,amplitude=4pt}] (-.5,-.65)--(2,-.65)node[midway,below=4pt,font=\scriptsize]{一个核苷酸的三部分};
\node[align=left,anchor=west,font=\scriptsize] at (8.9,.65) {糖和磷酸交替构成骨架；\\碱基连在糖上。};
\end{tikzpicture}}{核酸单链的组成示意（连线不表示所有具体化学键）\label{fig:11-nucleic}}

同一条链中，磷酸分别与相邻两个戊糖的羟基形成酯连接，得到\textbf{磷酸二酯键}。它与前面由羧酸形成的酯不是同一结构，但都可从“酸与醇羟基形成连接”来认识。DNA 的两条链之间，对应碱基通过氢键配对：A（腺嘌呤）与 T（胸腺嘧啶）配对，G（鸟嘌呤）与 C（胞嘧啶）配对。RNA 中以 U（尿嘧啶）代替 T；RNA 不应一律画成 DNA 那样的双螺旋。

\tu{\begin{tikzpicture}[font=\small]
\node at (0,0) {戊糖};\node at (1.2,0) {O};\node at (2.3,0) {P};\node at (3.4,0) {O};\node at (4.6,0) {戊糖};
\node at (2.3,1) {O};\node at (2.3,-1) {OH};
\draw (.45,0)--(.95,0) (1.4,0)--(2.1,0) (2.5,0)--(3.2,0) (3.65,0)--(4.15,0) (2.3,-.18)--(2.3,-.8);
\draw (2.25,.18)--(2.25,.8) (2.35,.18)--(2.35,.8);
\draw[decorate,decoration={brace,mirror,amplitude=4pt}] (.9,-1.4)--(3.7,-1.4)node[midway,below=4pt,font=\scriptsize]{相邻戊糖之间的磷酸二酯连接};
\node at (7,0.7) {A};\node at (10,0.7) {T};\node at (7,-.7) {G};\node at (10,-.7) {C};
\foreach \y in {.62,.78}{\draw[dashed] (7.5,\y)--(9.5,\y);}
\foreach \y in {-.9,-.7,-.5}{\draw[dashed] (7.5,\y)--(9.5,\y);}
\node[font=\scriptsize] at (8.5,1.2) {两条虚线：2个氢键};\node[font=\scriptsize] at (8.5,-1.3) {三条虚线：3个氢键};
\end{tikzpicture}}{链内磷酸二酯连接与DNA链间碱基配对（组成与连接示意，省略电离状态及碱基详细结构）\label{fig:11-dnapair}}

读图~\ref{fig:11-dnapair} 时，先区分实线与虚线：左图实线为共价连接，右图虚线为链间氢键。氢键使两链对应配对，同时允许在一定条件下分开；不能把它当成连接核苷酸的肽键。蛋白质、核酸、淀粉的单体分别为氨基酸、核苷酸、葡萄糖，结构不同，水解产物也不同。

\section{合成高分子：怎样连接，怎样拆回单体}

\subsection{单体、链节与聚合度}
聚乙烯由许多乙烯小分子连接而成。参加聚合的小分子叫\textbf{单体}，高分子链中重复出现的结构单元叫\textbf{链节}，重复次数 $n$ 叫\textbf{聚合度}。写作 \ce{-[CH2-CH2]_{$n$}-} 时，方括号两侧的键表示链还在继续，不能在每个碳上又补满成乙烷。材料中的分子通常长短不同，所以聚合度和相对分子质量一般取统计意义上的平均值。

\subsection{加聚：双键打开，侧基保留}
乙烯与溴加成时，双键中的一个 $\pi$ 键转变为新键；加聚也利用这个结构特点，只是新连接的对象换成了相邻的单体。原来双键的两个碳进入主链，每个碳原来所连的 H 或侧基仍留在原来的碳上。只从净变化看，不要求学习自由基或离子反应机理。

\tu{\begin{tikzpicture}[font=\small,>=Stealth]
\node at (0,0) {$n\,$};
\node (a) at (.65,0) {\ce{CH2}};\node (b) at (2,0) {CH};
\draw ([yshift=1.5pt]a.east)--([yshift=1.5pt]b.west);\draw[red!70!black,dashed] ([yshift=-1.5pt]a.east)--([yshift=-1.5pt]b.west);
\node (r) at (2,-.9) {R};\draw (b)--(r);
\draw[->] (3,0)--node[above]{聚合}(4.5,0);
\node (c) at (5.5,0) {\ce{CH2}};\node (d) at (6.9,0) {CH};\node (rr) at (6.9,-.9) {R};\draw(c)--(d)--(rr);
\draw[blue!70!black,thick] (4.95,0)--(c.west) (d.east)--(7.55,0);
\draw (5.05,.5)--(4.85,.5)--(4.85,-1.2)--(5.05,-1.2);
\draw (7.45,.5)--(7.65,.5)--(7.65,-1.2)--(7.45,-1.2);
\node at (7.9,-1.1) {$n$};
\node[align=left,anchor=west,font=\scriptsize] at (8.6,-.1) {红虚线：原双键中的一个键打开；\\蓝线：与相邻链节形成新连接。\\R 为 H、Cl、苯基等原有侧基。};
\end{tikzpicture}}{乙烯型单体的加聚：主链来自原双键碳\label{fig:11-addpoly}}

\[
\ce{$n$CH2=CHCl ->[\text{催化剂}] -[CH2-CHCl]_{$n$}-}
\]
加聚没有水等小分子副产物，链节的原子组成与单体相同。1,3-丁二烯可按 1,4-加聚生成以下高分子：
\[\ce{-[CH2-CH=CH-CH2]_{$n$}-}\]
链节中仍有一个双键，所以不能把“高分子”当作“没有不饱和键”。

\begin{liti}[推回乙烯型单体]
聚甲基丙烯酸甲酯（有机玻璃）的链节为 \ce{-CH2-C(CH3)(COOCH3)-}。写出单体，说明酯基在聚合中是否被消耗。
\tcblower
\jie 先找主链上的两个碳：\ce{-CH2-C-}；两个侧基 \ce{-CH3}、\ce{-COOCH3} 都连在第二个碳上。去掉链节两端与相邻单元的连接，在这两个碳之间恢复双键，得到 \ce{CH2=C(CH3)COOCH3}。聚合发生在碳碳双键，酯基作为侧基保留，不能把它拆成缩聚单体。
\end{liti}

\subsection{缩聚：官能团连接，同时生成小分子}
第 10 章的酯化把羧酸与醇连接一次；若分子上有两个或更多能反应的官能团，连接后两端仍有反应位置，就能继续接长。这个过程中同时生成水等小分子，叫\textbf{缩聚反应}。典型的两组配对是羧基与羟基形成酯基，羧基与氨基形成酰胺基。只含一个羟基的乙醇和只含一个羧基的乙酸只能连接成普通酯，不能仅靠这次反应不断长链。

\tu{\small
\begin{tabular}{c}
\ce{HO-R-COOH + HO-R-COOH -> HO-R-COO-R-COOH + H2O}\\[4pt]
\textcolor{blue!70!black}{中间形成酯基；左端仍有羟基，右端仍有羧基，可继续连接。}\\[5pt]
\ce{HO-R-COO-R-COOH ->[\text{继续连接}] H-[O-R-CO]_{$n$}-OH}\\
\end{tabular}}{羟基酸的缩聚过程：从两个分子看清为何能继续接长\label{fig:11-condpoly}}

数脱水量时先明确\textbf{链的形态和端基}。$n$ 个羟基酸分子连接成一条不成环的链，只需连接 $n-1$ 次，脱去 $(n-1)$ 个水分子；若首尾成环，还要再连接一次。$n$ 个二元酸和 $n$ 个二元醇共有 $2n$ 个单体分子，连接成一条直链要连接 $2n-1$ 次。带端基的完整方程式例如：
\[
\ce{$n$HOCH(CH3)COOH ->[\text{催化剂}][\text{△}] H-[OCH(CH3)CO]_{$n$}-OH + ($n$-1)H2O}
\]
\[
\begin{gathered}
 n\,\ce{HOOC-C6H4-COOH}+n\,\ce{HOCH2CH2OH}\\
 \xrightarrow{\text{催化剂、加热}}\ce{H-[OCH2CH2OCO-C6H4-CO]_{$n$}-OH}+(2n-1)\ce{H2O}
\end{gathered}
\]
后一产物是涤纶的主要成分（\ce{C6H4} 为对位亚苯基）。有些题只写内部链节、忽略端基，采用 $n$ 或 $2n$ 的简写计量；若要求完整方程式并写出端基，必须按上述具体结构守恒，不能把简写系数硬套进去。

\subsection{由聚酯、聚酰胺推回单体}
推单体就是将聚合过程倒过来看。不要一见方括号就“每两个碳切一次”，应先识别主链中的连接类型。

\begin{enumerate}
\item 主链为乙烯型的碳链：找原双键碳和保留的侧基，再恢复双键。主链有双键时，结合二烯烃加聚信息分析；共聚物还须按各单体的结构分段，不能任意二等分。
\item 主链有酯基 \ce{-CO-O-}：在羰基碳与单键氧之间断开，羰基碳一侧补 OH，氧一侧补 H，恢复羧酸和醇。
\item 主链有酰胺基 \ce{-CO-NH-}：断开羰基碳与 N 之间的键，羰基碳补 OH，N 补 H，恢复羧酸和胺。
\end{enumerate}

\tu{\small
\begin{tabular}{ccc}
聚酯连接处 & \ce{-R-CO}\textcolor{red!70!black}{$\,|\,$}\ce{O-R2-} & \ce{-R-COOH}\quad\ce{HO-R2-}\\[6pt]
聚酰胺连接处 & \ce{-R-CO}\textcolor{red!70!black}{$\,|\,$}\ce{NH-R2-} & \ce{-R-COOH}\quad\ce{H2N-R2-}\\[6pt]
&\textcolor{red!70!black}{竖线表示所断的 C—O 或 C—N 键}&补入的 OH 和 H 合计为一分子水\\
\end{tabular}}{用水解还原缩聚单体：先断连接，再补回水的两部分\label{fig:11-monomer}}

\begin{liti}[由单体检查回去]
尼龙-66 的内部链节为 \ce{-NH(CH2)6NHCO(CH2)4CO-}。推断其单体，并说明为何不是己二醇与己二酸。
\tcblower
\jie 链节主链含 \ce{-NH-CO-}，是酰胺连接而不是酯连接。把每个羰基碳与 N 的连接断开：\ce{-CO(CH2)4CO-} 两端各补 OH，得己二酸 \ce{HOOC(CH2)4COOH}；\ce{-NH(CH2)6NH-} 两端 N 各补 H，得己二胺 \ce{H2N(CH2)6NH2}。两者再缩合能恢复原链节。若用己二醇，主链将含 O，无法得到原来的 N 原子结构。
\end{liti}

\subsection{结构怎样影响材料性能}
线型高分子链之间主要靠分子间作用力结合，部分材料加热可软化、冷却又变硬，具有热塑性，如聚乙烯。若链之间通过共价键交联形成网状结构，受热不易重新流动成型，通常表现热固性，如充分固化的酚醛树脂。橡胶硫化是在链间引入含硫交联结构，可提高强度和弹性；交联键不必一律写成单一的二硫键。

涤纶、锦纶、腈纶分别属于聚酯、聚酰胺、聚丙烯腈纤维。高吸水性树脂既需要亲水基团与水相互作用，也需要适度交联以吸水膨胀而不直接溶散。离子交换树脂含可交换离子的基团，可用于水处理；导电高分子等材料则需依题给结构与性能信息分析，不必额外展开其微观机理。可降解聚酯的主链中有可水解的酯基；“来自天然原料”与“容易降解”没有简单的等价关系，应看具体结构和条件。

\section{有机合成：从一步转化走向路线设计}

\subsection{先确定哪些碳保留，再看官能团变化}
合成不是把所有反应条件试一遍。先给原料和目标物的碳编号，判断碳数、支链、环是否改变；若碳架不变，优先利用第 10 章的官能团转化。例如由乙烯合成乙酸：碳数一直为 2，只需经历烯烃加水、醇氧化、醛氧化。
\[
\ce{CH2=CH2 ->[\ce{H2O},\text{催化剂}] CH3CH2OH ->[\ce{O2},\ce{Cu},\text{△}] CH3CHO ->[\ce{O2},\text{催化剂}] CH3COOH}
\]
若目标碳数变多，就必须找形成 C—C 键的步骤。酯化虽能把两个有机分子接在一起，连接的是 C—O 键，不能说它增长了同一条连续碳链。高考常在题干提供增碳反应信息，重点是读懂新旧原子的对应关系，不必背诵大学有机合成反应名。

\begin{liti}[读取增碳信息]
已知 \ce{RCHO + HCN -> RCH(OH)CN}，腈基 \ce{-CN} 在酸性水溶液中水解可转化为 \ce{-COOH}。写出由乙醛制乳酸 \ce{CH3CH(OH)COOH} 的两步路线，并追踪新增碳来自哪里。
\tcblower
\jie 乙醛先加成：\ce{CH3CHO ->[\ce{HCN}] CH3CH(OH)CN}；再水解：\ce{CH3CH(OH)CN ->[\ce{H+},\ce{H2O}][\text{△}] CH3CH(OH)COOH}。原醛基碳成为连羟基的中间碳，HCN 的碳成为羧基碳，碳数由 2 增至 3。\ce{-CN} 连在碳上，不能误写成 \ce{-NC}。这是给定信息的应用，不表示 HCN 可供学生自行实验。
\end{liti}

缩短碳架的常见线索是强氧化造成的碳碳键断裂，例如丙烯的端位 \ce{CH2=} 可变为 \ce{CO2}，另一段变为乙酸；苯的同系物侧链在符合条件时氧化为羧基。成环则要找分子两端可相互反应的官能团，例如羟基酸的分子内酯化。具体断裂产物或成环规律若超出课本，必须有题给信息支撑。

\subsection{从目标物逆推最后一步}
\textbf{逆合成分析}是先问“目标物最后可由哪一步反应得到”，再对所需前体重复这个问题，直到回到给定原料。逆推箭头 $\Rightarrow$ 只表示分析方向，正式路线仍应按实际反应方向书写，并写全试剂和条件。

\begin{liti}[从一个酯开始逆推]
以乙醇为唯一有机原料，无机试剂任选，合成乙酸乙酯。写出分析过程和正式路线。
\tcblower
\silu 目标含酯基，最后一步可由乙酸与乙醇酯化；乙醇已有，缺乙酸。乙酸可由乙醛氧化，乙醛可由乙醇氧化。\par
\jie 将原料分为两份：一份保留作醇，另一份沿“乙醇$\to$乙醛$\to$乙酸”制酸。最后反应：
\[\ce{CH3COOH + CH3CH2OH <=>[\text{浓硫酸}][\text{△}] CH3COOCH2CH3 + H2O}\]
前两步分别用 \ce{O2}/Cu、加热和 \ce{O2}/催化剂。题目说“唯一有机原料”，允许从它制得多种中间体再汇合，不能直接另取乙酸作原料。
\end{liti}

\subsection{反应顺序与官能团保护}
多官能团分子可能有两个位置都与同一试剂反应。优先选择合适试剂实现\textbf{选择性转化}：例如把 \ce{CH2=CHCHO} 的醛基氧化而保留碳碳双键，可选银氨溶液或新制 \ce{Cu(OH)2}，反应后酸化；酸性 \ce{KMnO4} 会同时破坏双键，不能用来完成这一目标。这里没有先改造再恢复双键，不能称为“保护后脱保护”。

只有选择性不足时，才考虑\textbf{暂时把某官能团变为不参与目标反应的结构，主反应后再恢复}。例如题目若给出“酚羟基可转化为较耐氧化的醚、反应后再恢复酚羟基”的信息，才能用来设计氧化侧链而保留酚羟基的方案。保护反应、主体反应、脱保护都要互相兼容；不能笼统地用 HBr 加成来保护任意双键，因为它可能同时使醇羟基发生取代，消去恢复时还可能产生位置异构体。

\begin{liti}[由信息安排顺序]
目标是把 \ce{HOCH2CH2CHO} 中的羟基氧化为羧基，同时最终保留原来的醛基。题给信息：①醛基可与乙二醇形成缩醛，在碱性氧化条件下稳定；②缩醛在酸性水溶液中可水解恢复醛基；③题给碱性氧化剂能把伯醇氧化为羧酸盐。写出三步操作的先后次序，不必写缩醛具体结构。
\tcblower
\jie 先用乙二醇将原醛基转化为缩醛；再用题给氧化剂氧化原来的 \ce{-CH2OH}，得到羧酸盐；最后酸化、水解缩醛，恢复醛基，得到 \ce{HOOCCH2CHO}。若先氧化，醛基更易被氧化；若用酸性氧化条件，保护基可能先水解。可行性来自三条题给信息的共同限制。
\end{liti}

\subsection{把新反应接入熟悉的路线}
陌生信息先读“谁的哪根键变化”，再读碳数与官能团是否守恒。以下综合例题不要求记忆新反应机理，只要求将一个已给出的增碳反应接入熟悉的醇、醛、酸、酯转化。

\begin{liti}[逐步搭建综合路线]
已知 \ce{RCHO + CH3CHO ->[\text{稀 NaOH}] RCH(OH)CH2CHO ->[\text{△}] RCH=CHCHO + H2O}。以乙醇为唯一有机原料，合成 \ce{CH3CH=CHCOOCH2CH3}，无机试剂任选。说明每步目的。
\tcblower
\jie 目标是酯，逆推为 2-丁烯酸 \ce{CH3CH=CHCOOH} 与乙醇。保留一份乙醇；另一份催化氧化为乙醛：\ce{2CH3CH2OH + O2 ->[\ce{Cu}][\text{△}] 2CH3CHO + 2H2O}。两个乙醛分子按题给信息连接，先得 \ce{CH3CH(OH)CH2CHO}，再脱水得 \ce{CH3CH=CHCHO}，碳数由 2 变为 4。用碱性新制 \ce{Cu(OH)2} 氧化醛基，酸化得 \ce{CH3CH=CHCOOH}，保留双键；最后与保留的乙醇在浓硫酸、加热条件下酯化。每步均能指出变化的官能团，碳数也与目标一致。
\end{liti}

\zhangmo{本章小结}
\subsection*{知识总结}
糖类先按水解产物分类，再按结构判断还原性。油脂看三个酯基及烃基中的双键；蛋白质看氨基酸之间的肽键及整条链的空间结构；核酸看核苷酸及糖—磷酸骨架。它们不能只凭“营养物质”这一共同用途混在一起记忆。

高分子题先分清加聚与缩聚，再确定主链、侧基和连接处。加聚把原双键碳接入主链；缩聚利用官能团反应，推单体可借水解把连接断开。写完整方程式要检查端基、原子守恒和小分子数。

\subsection*{方法提升}
合成路线先核对碳骨架，再核对官能团位置；从目标逆推缺少的前体，按正向路线检查条件兼容。聚合计量从实际连接次数出发，所得单体再正向连接一次核查。
\subsection*{易错点归纳}
常见失误是把酯化误认为增长连续碳链、用强氧化剂破坏应保留的双键、漏写水解后的酸化，以及把题给的新反应当作无需条件的普遍规律。

\zhangmo{分层练习}
\zu{A组\quad 基础巩固}
\begin{xiti}[start=1]
\item \laiyuan{基础练习} 下列说法正确的是\xz
\xxa{淀粉和纤维素通式相同，互为同分异构体}{蔗糖与麦芽糖完全水解的产物均只有葡萄糖}{向蛋白质溶液中加入浓硫酸铵溶液析出沉淀，加水后又溶解，属于盐析}{油脂与蛋白质都是高分子，完全水解产物都含氨基酸}
\item \laiyuan{本书编写} 蔗糖与麦芽糖各 0.10 mol，分别完全水解。所得葡萄糖的物质的量各为多少？各自产物的总质量是否相同？
\item \laiyuan{本书编写} 写出甘氨酸分别与 NaOH、HCl 反应的产物。10 个氨基酸分子连接成一条不成环的直链多肽，形成多少个肽键、脱去多少个水分子？
\item \laiyuan{基础练习} 聚乳酸的内部链节为 \ce{-O-CH(CH3)-CO-}，写出单体和反应类型。充分水解后得到的产物是否属于糖？说明结构依据。
\item \laiyuan{本书编写} 加热使蛋白质凝固，与用酸使其充分水解，是否是同一个过程？DNA 中链内骨架和两条链之间分别靠哪类相互作用连接？某双链 DNA 共含 100 对碱基，其中 A 共 40 个，求 G 与 C 的总数及碱基对间氢键总数。
\end{xiti}

\zu{B组\quad 高考典型}
\begin{xiti}
\item \laiyuan{本书编写} 为检验淀粉的酸性水解液，甲取一份直接加入新制 \ce{Cu(OH)2} 并加热，无砖红色沉淀，便判断未水解；乙另取一份加碘水，不变蓝，便判断已完全水解。分别指出两种判断的不足，给出完整方案。
\item \laiyuan{基础练习} 写出以下高分子的单体：\par (1) 涤纶 \ce{H-[OCH2CH2OCO-C6H4-CO]_{$n$}-OH}，\ce{C6H4} 为对位亚苯基；\par (2) \ce{-[CH2-C(CH3)(COOCH3)]_{$n$}-}；\par (3) 丁苯橡胶的一个理想化内部结构单元 \ce{-CH2-CH=CH-CH2-CH2-CH(C6H5)-}。第三问结合二烯烃的 1,4-加聚分析。
\item \laiyuan{本书编写} $n$ 个乳酸分子全部连接成一条直链，所得链有几个酯基、脱去几个水分子？若同样 $n$ 个分子首尾连接成环，二者有何变化？分别写出产物分子式（用 $n$ 表示）。
\item \laiyuan{基础练习} 阿司匹林可写为 \ce{CH3COO-C6H4-COOH}（两个取代基在邻位）。(1) 写出含氧官能团名称。(2) 1 mol 阿司匹林与足量 NaOH 溶液共热，最多消耗多少摩尔 NaOH？写出方程式。(3) 由水杨酸与乙酸酐 \ce{(CH3CO)2O} 反应制取阿司匹林，同时生成乙酸，写出反应关系并说明反应位置。
\end{xiti}

\zu{C组\quad 冲刺高分}
\begin{xiti}
\item \laiyuan{综合练习} 肉桂酸乙酯为 \ce{C6H5CH=CHCOOC2H5}。已知以下两步反应：
\[\ce{RCHO + CH3CHO ->[\text{稀 NaOH}] RCH(OH)CH2CHO}\]
\[\ce{RCH(OH)CH2CHO ->[\text{△}] RCH=CHCHO + H2O}\]
现有路线：甲苯 $\xrightarrow{\ce{Cl2},\text{光照}}$ B $\xrightarrow{\text{NaOH 水溶液，△}}$ C $\xrightarrow{\ce{O2},\ce{Cu},\text{△}}$ D $\xrightarrow{\text{与乙醛按题给反应}}$ E（\ce{C9H8O}）$\xrightarrow{\text{氧化、酸化}}$ F $\xrightarrow{\text{与乙醇酯化}}$ 产品。(1) 写出 B～F 的结构简式。(2) E$\to$F 应选哪类氧化剂？为什么？(3) F 的同分异构体中，同时含苯环、羧基、碳碳双键，且苯环上恰有两个取代基的有几种（不考虑立体异构）？展开计数依据。
\item \laiyuan{本书编写} 以乙醇为唯一有机原料、无机试剂任选，根据本章给出的 \ce{HCN} 加成和腈基水解信息，设计合成聚乳酸的路线。(1) 写出每步中间体与条件。(2) 若起始有 0.50 mol 乙醇，各步定量转化，最终只形成一条含 $n$ 个乳酸残基的直链（$n$ 很大），忽略端基对质量的影响，理论可得聚乳酸多少克？(3) 为什么这条路线并非所有碳原子都来自乙醇？
\end{xiti}

\zhangmo{练习解析}
\begin{jiexilist}
\item \ans{C}淀粉和纤维素的聚合度不确定，不能只凭通式认定分子式相同；蔗糖还生成果糖；浓硫酸铵引起的可逆析出是盐析；油脂不是高分子，水解得甘油和高级脂肪酸（或盐），蛋白质水解得氨基酸。
\item 蔗糖生成葡萄糖 0.10 mol、果糖 0.10 mol；麦芽糖生成葡萄糖 0.20 mol。两者都吸收 0.10 mol 水，产物总质量均为 $0.20\times180=36.0$ g。产物总质量相同不表示产物种类相同。
\item 与 NaOH 生成 \ce{H2NCH2COONa} 和水；与 HCl 生成 \ce{[H3NCH2COOH]Cl}。10 个单体连成不成环的一条链要连接 9 次，形成 9 个肽键、脱去 9 个水分子。若分子还有可反应侧链或形成环，须重新按实际连接数计算。
\item 单体 \ce{CH3CH(OH)COOH}（乳酸），缩聚反应。乳酸是羟基羧酸，不是多羟基醛或多羟基酮，不属于糖；含羟基不能作为判糖的充分条件。
\item 不是。变性主要破坏或改变维持原有生理功能所需的结构，不等于肽键全部断裂；充分水解则断开肽键生成小分子。DNA 同一条链的糖—磷酸骨架靠磷酸二酯键等共价连接，互补碱基之间靠氢键配对。共100对即200个碱基，A与T各40个，所以G、C各60个，共120个；A—T为40对，G—C为60对，氢键总数为 $40\times2+60\times3=260$ 个。
\item 甲未排除酸的干扰：酸会破坏检验所需的碱性，阴性不能证明无葡萄糖。乙只有“不变蓝”的单一现象，须确认碘试剂和原样有效，并检验水解产物。正确做法是分两份：一份冷却后加碘水；另一份先用 NaOH 调至碱性，再加入新制 \ce{Cu(OH)2} 并加热。既变蓝又有砖红沉淀为部分水解；不变蓝而有砖红沉淀为完全水解；变蓝且无砖红沉淀，在操作有效时判断未水解。
\item (1) 对苯二甲酸 \ce{HOOC-C6H4-COOH} 和乙二醇 \ce{HOCH2CH2OH}。把主链酯基的 C—O 键断开，酸端补 OH、醇端补 H 即得。(2) \ce{CH2=C(CH3)COOCH3}，酯基为侧基，主链两个碳恢复双键。(3) \ce{CH2=CH-CH=CH2} 和 \ce{CH2=CHC6H5}。先把以主链双键为中心的四个碳还原为丁二烯，再将余下两个相邻主链碳恢复为苯乙烯。
\item 直链有 $n-1$ 个酯基，脱去 $n-1$ 个水；分子式为 \ce{C_{$3n$}H_{$4n+2$}O_{$2n+1$}}，由 $n\ce{C3H6O3}-(n-1)\ce{H2O}$ 得到。成环多连接一次，有 $n$ 个酯基、脱去 $n$ 个水，分子式为 \ce{C_{$3n$}H_{$4n$}O_{$2n$}}。端基的差别正好对应一个水分子。
\item (1) 羧基、酯基。(2) 3 mol：原羧基消耗 1 mol，酚酯水解先消耗 1 mol，水解生成的酚羟基再消耗 1 mol。
\[\ce{CH3COO-C6H4-COOH + 3NaOH ->[\text{△}] NaO-C6H4-COONa + CH3COONa + 2H2O}\]
(3) \ce{HO-C6H4-COOH + (CH3CO)2O -> CH3COO-C6H4-COOH + CH3COOH}。酚羟基上的 H 被乙酰基代替，羧基保留；属于取代反应。
\item (1) B \ce{C6H5CH2Cl}，C \ce{C6H5CH2OH}，D \ce{C6H5CHO}，E \ce{C6H5CH=CHCHO}，F \ce{C6H5CH=CHCOOH}。光照取代在甲苯侧链；D 与乙醛按信息增碳、脱水。(2) 选银氨溶液或碱性新制 \ce{Cu(OH)2}，反应后酸化，氧化醛基而保留双键；酸性 \ce{KMnO4} 也会氧化断裂双键。(3) F 为 \ce{C9H8O2}，苯环、羧基和双键已用完所有不饱和度。除苯环外仅有 3 个碳，其中羧基占 1 个，双键占 2 个，因此两个取代基只能为 \ce{-COOH} 与 \ce{-CH=CH2}，邻、间、对共 3 种，不能再加入甲基或 \ce{CH2}。
\item (1) \ce{CH3CH2OH ->[\ce{O2},\ce{Cu}][\text{△}] CH3CHO ->[\ce{HCN}] CH3CH(OH)CN ->[\ce{H+},\ce{H2O}][\text{△}] CH3CH(OH)COOH}，乳酸在催化剂、加热条件下缩聚为 \ce{H-[OCH(CH3)CO]_{$n$}-OH}。(2) 0.50 mol 乙醇对应 0.50 mol 乳酸残基；链节 \ce{C3H4O2} 的式量为 72，质量约 $0.50\times72=36.0$ g。精确处理一条链的两个端基时，聚合物总组成比所有内部链节的组成多一个水，但本问已忽略这一影响。(3) 每个残基有 3 个碳，乙醇只提供其中 2 个；另 1 个来自无机原料 HCN 的碳，最后成为羰基碳。“唯一有机原料”不等于“唯一含碳原料”。
\end{jiexilist}
